%%% ============================================================
%%% Joint Closure of Two Commutative Binary Operations on Z/10Z:
%%% an Eight-Element Chain and a Normalizer Identity
%%%
%%% B.R. Sanders, M. Gish (2026)
%%% Target venue: Algebraic Combinatorics
%%% Backup: Communications in Algebra; Discrete Mathematics
%%% MSC: 20N02, 20N05, 11A07
%%%
%%% This is the SEED-NARROW extraction from the bundled
%%% four_core_FINAL_BUNDLED draft (review round 3).  Theorems
%%% 1 and 2 (chain + normalizer) only.  Theorems 3 (closed-form
%%% attractor), 4 (Galois D_4), and 5 (alpha-sweep) are deferred
%%% to companion papers in preparation, listed in the Forward
%%% Directions section.
%%%
%%% Verification: 4core_verification.py archived at DOI
%%% 10.5281/zenodo.18852047 (chain enumeration + normalizer
%%% symbolic checks; the attractor / Galois / alpha-sweep checks
%%% in the same script support the companion papers in
%%% preparation).
%%% ============================================================

\documentclass[11pt,reqno]{amsart}

\usepackage{amsmath,amssymb,amsthm,mathtools}
\usepackage[margin=1in]{geometry}
\usepackage[colorlinks=true,linkcolor=blue,citecolor=blue,urlcolor=blue]{hyperref}
\usepackage{microtype}
\usepackage{array}

%% -------- theorem environments --------
\theoremstyle{plain}
\newtheorem{theorem}{Theorem}
\newtheorem{lemma}[theorem]{Lemma}
\newtheorem{proposition}[theorem]{Proposition}
\newtheorem{corollary}[theorem]{Corollary}

\theoremstyle{definition}
\newtheorem{definition}[theorem]{Definition}

\theoremstyle{remark}
\newtheorem*{remark}{Remark}
\newtheorem*{notation}{Notation}

%% -------- macros --------
\newcommand{\Z}{\mathbb{Z}}
\newcommand{\Q}{\mathbb{Q}}
\newcommand{\R}{\mathbb{R}}
\newcommand{\Zten}{\Z/10\Z}
\newcommand{\Cfour}{\mathcal{C}}
\newcommand{\Tfuse}{\widehat{T}}
\newcommand{\Bfuse}{\widehat{B}}

\title[Joint closure on $\Zten$]{Joint Closure of Two Commutative Binary
Operations on $\Zten$:\\ an Eight-Element Chain and a Normalizer Identity}

\author{B.~R.~Sanders}
\address{7Site LLC, Hot Springs, AR, USA}
\email{brayden@7site.io}

\author{M.~Gish}
\address{Independent Researcher}

\subjclass[2020]{Primary 20N02; Secondary 20N05, 11A07}

\keywords{commutative magma, sub-magma chain, joint closure,
fuse normalizer, polynomial identity}

\date{\today}

\begin{document}

\begin{abstract}
We study a pair of commutative binary operations
$T,B:\Zten\times\Zten\to\Zten$ given explicitly by their Cayley
tables, both of which arise as restrictions of a common
operator-substrate construction (detailed in the companion
paper~\cite{SandersGish2026Sigma}). Two exact structural results
emerge from direct enumeration and elementary algebra.

(i) The sub-magmas of $\Zten$ that are jointly closed under
both $T$ and $B$ form a strict $8$-element chain in the inclusion
order, with sizes $\{1, 4, 5, 6, 7, 8, 9, 10\}$; sizes $2$ and $3$
are forbidden as joint closures. Among all $1023$ non-empty
subsets of $\Zten$, exactly these eight satisfy joint closure.

(ii) On the minimal non-trivial element of this chain, the
four-element subset $\Cfour=\{0,7,8,9\}$, the two operations have
a common fuse-normalizer
$Z_{T}(p)=Z_{B}(p)=(p_{0}+p_{7}+p_{8}+p_{9})^{2}$ for every
$p\in\R^{10}$ supported on $\Cfour$, despite the fact that $T$
and $B$ disagree on $12$ of the $16$ cells of the $4\times 4$
restriction. The identity is a polynomial identity in the four
$\Cfour$-coordinates and is independently verifiable by symbolic
expansion.

The chain admits a clean structural reading in terms of a
specific permutation $\sigma:\Zten\to\Zten$: the chain walks the
$\sigma$-forward orbit of $7$, with the $\sigma$-fixed elements
of $\{0,3,8,9\}$ contributing at the bottom, the size-$1\to 4$
jump, and the size-$7\to 8$ bridge step.

The dynamical and number-theoretic phenomena that arise from
$(T,B)$ --- a closed-form fixed point of the convex-combination
iteration at $\alpha=\tfrac12$ with $p_{7}/p_{8}=1+\sqrt{3}$, an
associated quartic Galois extension, and observed
mixing-weight phenomena at $\alpha\in(0,1)$ --- are the subject
of two companion papers in preparation
(see Forward Directions, \S\ref{sec:forward}).
\end{abstract}

\maketitle

%% -------- section 1: introduction --------
\section{Introduction}\label{sec:intro}

This paper studies a small object: the family of sub-magmas of
$\Zten$ that are simultaneously closed under two specific
commutative binary operations $T$ and $B$. We establish two
exact structural results, accompanied by a deterministic
verification script.

The two operations are presented as Cayley tables in
\S\ref{sec:setup}. Both are commutative and both have $0$
acting as a left zero on the off-diagonal. Beyond these features
they differ markedly: $T$ has the value $7$ acting as a row
absorber on most of $\Zten$; $B$ is approximately incrementing
in the sense that $B(j,j) = j+1 \pmod{10}$ for
$j \in \{1,\dots,7\}$, with diagonal exceptions at
$j \in \{0,8,9\}$ (where $B(0,0)=0$, $B(8,8)=7$, and $B(9,9)=0$).
Despite this surface dissimilarity, two exact structural
phenomena link them tightly: a complete classification of the
joint-closure lattice into a single $8$-element chain
(Theorem~\ref{thm:chain}), and a polynomial identity for the
fuse-normalizer of the two operations on the chain's minimal
non-trivial element (Theorem~\ref{thm:norm}).

\paragraph{The chain (Theorem~\ref{thm:chain}).} The sub-magmas of
$\Zten$ jointly closed under both $T$ and $B$, partially ordered
by inclusion, form a totally ordered $8$-element chain
\[
\{0\}\subset\Cfour\subset\Cfour\cup\{6\}\subset\Cfour\cup\{5,6\}\subset
\Cfour\cup\{4,5,6\}\subset\Cfour\cup\{3,4,5,6\}\subset
\Cfour\cup\{2,3,4,5,6\}\subset\Zten,
\]
where $\Cfour=\{0,7,8,9\}$. The chain is strict: no size-$2$ or
size-$3$ subset of $\Zten$ is jointly closed. Among all $1023$
non-empty subsets of $\Zten$, exactly these eight satisfy joint
closure.

\paragraph{The 4-core normalizer identity (Theorem~\ref{thm:norm}).}
On the minimal non-trivial element $\Cfour$ of the chain, the
fuse-normalizer
$Z_{M}(p) = \sum_{c\in\Zten}\sum_{(i,j):M(i,j)=c}p_{i}p_{j}$ is
identical for $M = T$ and $M = B$:
\[
Z_{T}(p) = Z_{B}(p) = (p_{0}+p_{7}+p_{8}+p_{9})^{2}\qquad\text{for }p\text{ supported on }\Cfour,
\]
even though $T$ and $B$ disagree on $12$ of the $16$ cells of the
$4\times 4$ restriction. The identity is a polynomial identity
in the four $\Cfour$-coordinates, exact and independently
verifiable by symbolic expansion.

\paragraph{The pair $(T,B)$ is not independent.} Both Cayley
tables arise from a common operator-substrate construction in
which $T$ is the rank-3 threshold projection at threshold
$T^{*} = 5/7$ and $B$ is the full rank-10 composition; the
construction is detailed in the companion
paper~\cite{SandersGish2026Sigma}. The present paper takes
$(T, B)$ at $N = 10$ as given and studies the joint-closure
structure of the pairing. The chain rigidity above has been
verified to persist over $\mathbb{F}_{p}$-lifts of the same
table for $p \in \{2, 3, 5, 7, 11, 13\}$; the
$\mathbb{F}_{p}$-universality of the joint-closure structure is
the subject of a companion paper in preparation.

\paragraph{Positioning.} The classical Drápal--Kepka tradition
of low-density associativity in
quasigroups~\cite{Kepka1980, DrapalKepka1985} treats
non-associativity counts on cancellative structures. The present
paper studies the joint closure of two non-cancellative magmas
and the polynomial identity arising from cell-cancellation
across the $4\times 4$ restriction; it is not implied by the
corresponding density results in those references.

\paragraph{Outline.} \S\ref{sec:setup} fixes the two tables.
\S\ref{sec:chain} proves the joint-closure chain
(Theorem~\ref{thm:chain}). \S\ref{sec:4core} records the
corollary that $\Cfour=\{0,7,8,9\}$ is jointly closed and is the
minimal non-trivial element of the chain. \S\ref{sec:norm}
proves the normalizer coincidence (Theorem~\ref{thm:norm}).
\S\ref{sec:verify} catalogs the verification script.
\S\ref{sec:scope} states the limitations of the present results.
\S\ref{sec:forward} lists three companion papers in preparation
that develop the dynamical and number-theoretic structure of the
pair $(T,B)$.

%% -------- section 2: setup --------
\section{Definitions and the two tables}\label{sec:setup}

We work throughout on the additive group $\Zten$, identified with
the labels $\{0,1,\dots,9\}$.

\begin{definition}[The two tables]\label{def:tables}
The two binary operations $T,B:\Zten\times\Zten\to\Zten$ are given
by their Cayley tables:
\[
T = \left[
\begin{array}{*{10}{c}}
0 & 0 & 0 & 0 & 0 & 0 & 0 & 7 & 0 & 0 \\
0 & 7 & 3 & 7 & 7 & 7 & 7 & 7 & 7 & 7 \\
0 & 3 & 7 & 7 & 4 & 7 & 7 & 7 & 7 & 9 \\
0 & 7 & 7 & 7 & 7 & 7 & 7 & 7 & 7 & 3 \\
0 & 7 & 4 & 7 & 7 & 7 & 7 & 7 & 8 & 7 \\
0 & 7 & 7 & 7 & 7 & 7 & 7 & 7 & 7 & 7 \\
0 & 7 & 7 & 7 & 7 & 7 & 7 & 7 & 7 & 7 \\
7 & 7 & 7 & 7 & 7 & 7 & 7 & 7 & 7 & 7 \\
0 & 7 & 7 & 7 & 8 & 7 & 7 & 7 & 7 & 7 \\
0 & 7 & 9 & 3 & 7 & 7 & 7 & 7 & 7 & 7
\end{array}
\right],\qquad
B = \left[
\begin{array}{*{10}{c}}
0 & 1 & 2 & 3 & 4 & 5 & 6 & 7 & 8 & 9 \\
1 & 2 & 3 & 4 & 5 & 6 & 7 & 2 & 6 & 6 \\
2 & 3 & 3 & 4 & 5 & 6 & 7 & 3 & 6 & 6 \\
3 & 4 & 4 & 4 & 5 & 6 & 7 & 4 & 6 & 6 \\
4 & 5 & 5 & 5 & 5 & 6 & 7 & 5 & 7 & 7 \\
5 & 6 & 6 & 6 & 6 & 6 & 7 & 6 & 7 & 7 \\
6 & 7 & 7 & 7 & 7 & 7 & 7 & 7 & 7 & 7 \\
7 & 2 & 3 & 4 & 5 & 6 & 7 & 8 & 9 & 0 \\
8 & 6 & 6 & 6 & 7 & 7 & 7 & 9 & 7 & 8 \\
9 & 6 & 6 & 6 & 7 & 7 & 7 & 0 & 8 & 0
\end{array}
\right].
\]
Both $T$ and $B$ are commutative.
\end{definition}

\begin{notation}
We refer to the indices $0,7,8,9$ frequently. In proofs we
sometimes write $(p_{0},p_{7},p_{8},p_{9})=(v,h,\beta,r)$ for
masses on these four indices.
\end{notation}

\begin{definition}[Sub-magma; joint closure]\label{def:submagma}
A subset $S\subseteq\Zten$ is a \emph{sub-magma} of $(\Zten,M)$ if
$M(i,j)\in S$ for all $i,j\in S$. We call $S$ \emph{jointly
closed} (under $T$ and $B$) if it is a sub-magma of both. The set
of jointly closed subsets, ordered by inclusion, is denoted
$\mathcal{L}(T,B)$.
\end{definition}

\begin{definition}[Fuse and normalizer]\label{def:fuse}
For a binary operation $M:\Zten\times\Zten\to\Zten$ and a vector
$p=(p_{0},\dots,p_{9})\in\R^{10}$, the \emph{fuse}
$\widehat{M}(p)\in\R^{10}$ is the convolution along the operation:
\[
\widehat{M}(p)_{c}\;=\;\sum_{(i,j)\in\Zten^{2}\,:\,M(i,j)=c}p_{i}p_{j},
\qquad c\in\Zten.
\]
The \emph{normalizer} is $Z_{M}(p)=\sum_{c\in\Zten}\widehat{M}(p)_{c}$.
\end{definition}

\begin{remark}
By bilinearity, $Z_{M}(p) = (\sum_{i}p_{i})^{2}$ identically for any
$M$ and any $p$, since each ordered pair $(i,j) \in \Zten^{2}$
contributes $p_{i}p_{j}$ to exactly one output bucket. The
non-trivial content of Theorem~\ref{thm:norm} below is the
identity \emph{restricted to $p$ supported on $\Cfour$}: there,
the fuses $\Tfuse(p)$ and $\Bfuse(p)$ have support contained in
$\Cfour$ (since $\Cfour$ is jointly closed,
Corollary~\ref{cor:4core}), so the four output coordinates
$\widehat{M}(p)_{0}, \widehat{M}(p)_{7}, \widehat{M}(p)_{8},
\widehat{M}(p)_{9}$ exhaust $Z_{M}(p)$. The identity
$Z_{T}(p) = Z_{B}(p) = (v+h+\beta+r)^{2}$ on $\Cfour$-supported
$p$ is then a four-coordinate polynomial identity, distinct from
the trivial total-mass identity, and arises from a non-trivial
cell-by-cell cancellation across the $12$ of the $16$ cells on
which $T$ and $B$ disagree.
\end{remark}

%% -------- section 3: chain --------
\section{The joint-closure chain}\label{sec:chain}

\begin{theorem}[Joint-closure chain]\label{thm:chain}
The sub-magmas of $\Zten$ jointly closed under both $T$ and $B$
form a strict chain of length $8$ in inclusion order:
\[
\mathcal{L}(T,B)\;=\;\bigl\{S_{1}\subset S_{4}\subset S_{5}\subset
S_{6}\subset S_{7}\subset S_{8}\subset S_{9}\subset S_{10}\bigr\},
\]
where
\begin{align*}
S_{1}&=\{0\}, & S_{4}&=\{0,7,8,9\}, \\
S_{5}&=\{0,6,7,8,9\}, & S_{6}&=\{0,5,6,7,8,9\}, \\
S_{7}&=\{0,4,5,6,7,8,9\}, & S_{8}&=\{0,3,4,5,6,7,8,9\}, \\
S_{9}&=\{0,2,3,4,5,6,7,8,9\}, & S_{10}&=\Zten.
\end{align*}
The chain is strict: no subset of $\Zten$ of size $2$ or size $3$
is jointly closed.
\end{theorem}

\begin{proof}
We verify the joint-closure status of each candidate subset by
direct computation. The proof has two parts: (a) confirming that
each $S_{k}$ in the displayed list is jointly closed, and (b)
confirming that no other non-empty subset is jointly closed.

\textbf{Part (a) --- the eight subsets are jointly closed.} For
each $k\in\{1,4,5,6,7,8,9,10\}$, the verification reduces to
checking that $T(i,j)\in S_{k}$ and $B(i,j)\in S_{k}$ for all
$i,j\in S_{k}$. We display the size-$4$ case explicitly and refer
to the verification script (\S\ref{sec:verify}) for the others.

For $S_{4}=\{0,7,8,9\}$, the $4\times 4$ restrictions are
\[
T\big|_{S_{4}\times S_{4}}=\begin{pmatrix}
0 & 7 & 0 & 0 \\
7 & 7 & 7 & 7 \\
0 & 7 & 7 & 7 \\
0 & 7 & 7 & 7
\end{pmatrix},
\qquad
B\big|_{S_{4}\times S_{4}}=\begin{pmatrix}
0 & 7 & 8 & 9 \\
7 & 8 & 9 & 0 \\
8 & 9 & 7 & 8 \\
9 & 0 & 8 & 0
\end{pmatrix}.
\]
Every entry of $T\big|_{S_{4}\times S_{4}}$ lies in
$\{0,7\}\subseteq S_{4}$. Every entry of
$B\big|_{S_{4}\times S_{4}}$ lies in $\{0,7,8,9\}=S_{4}$.

\textbf{Part (b) --- no other subset is jointly closed.}

\emph{Sizes $2$ and $3$.} We dispatch these by exhaustive case
analysis on the diagonal entries of $B$. Inspecting the diagonal:
\[
B(j,j)=\begin{cases}
0 & j=0,\\
j+1 & j\in\{1,2,3,4,5,6,7\},\\
7 & j=8,\\
0 & j=9.
\end{cases}
\]

\emph{Size $2$.} Let $S=\{a,b\}$ with $a<b$. Joint closure requires
$B(a,a),B(b,b)\in S$ and $T(a,a),T(b,b)\in S$, plus the off-diagonal
constraints.

\emph{Subcase $0\in S$.} Then $S=\{0,a\}$ with $a\in\{1,\dots,9\}$.
The diagonal $B(a,a)$ must lie in $\{0,a\}$. From the diagonal
table: for $a\in\{1,\dots,7\}$, $B(a,a)=a+1\notin\{0,a\}$; for
$a=8$, $B(8,8)=7\notin\{0,8\}$; for $a=9$, $B(9,9)=0\in\{0,9\}$ and
$B(0,9)=B(9,0)=9\in\{0,9\}$, so $\{0,9\}$ is $B$-closed. However
$T(9,9)=7\notin\{0,9\}$, so $\{0,9\}$ fails $T$-closure.

\emph{Subcase $0\notin S$.} Then $S\subseteq\{1,\dots,9\}$, and
both $B(a,a)\in S$ and $B(b,b)\in S$ are required. Since
$B(j,j)\neq j$ for all $j\in\{1,\dots,9\}$, we need $B(a,a)=b$ and
$B(b,b)=a$. Listing all $j\in\{1,\dots,9\}$ with their diagonal
values: $(1,2),(2,3),(3,4),(4,5),(5,6),(6,7),(7,8),(8,7),(9,0)$.
The only unordered pair $\{a,b\}\subseteq\{1,\dots,9\}$ with
$B(a,a)=b$ and $B(b,b)=a$ is $\{7,8\}$ (since $B(7,7)=8$ and
$B(8,8)=7$). However $B(7,8)=9\notin\{7,8\}$, so $\{7,8\}$ fails
$B$-closure.

The remaining size-$2$ candidates do not satisfy
$B(a,a),B(b,b)\in S$ and are excluded immediately. No size-$2$
subset is jointly closed.

\emph{Size $3$.} Suppose $S=\{a,b,c\}$ with $a<b<c$ is jointly
closed. The three diagonal entries $B(a,a),B(b,b),B(c,c)$ and the
three diagonals $T(s,s)$ for $s\in S$ must all lie in $S$. Direct
enumeration of $\binom{10}{3}=120$ candidates against just the
six diagonal constraints leaves only two candidates:
$\{0,7,8\}$ and $\{6,7,8\}$. (For $\{0,7,8\}$: $B$-diagonals are
$0,8,7$; for $\{6,7,8\}$: $B$-diagonals are $7,8,7$; in both
cases the $T$-diagonals are $0,7,7$ and $7,7,7$ respectively,
all in $S$.) However both candidates fail $B$-closure at the
off-diagonal cell $B(7,8)=9$: in $\{0,7,8\}$, $9\notin S$; in
$\{6,7,8\}$, $9\notin S$. Hence no size-$3$ subset is jointly
closed.

\emph{Sizes $4$ through $10$ --- uniqueness within size class.}
We must confirm that the listed $S_{k}$ is the \emph{unique}
jointly-closed subset of its size class, for $k\in\{4,5,6,7,8,9\}$.
A useful structural reduction: \emph{any jointly-closed subset
$S \subseteq \Zten$ of size at least $4$ contains $\Cfour =
\{0, 7, 8, 9\}$.} Indeed, suppose $|S| \geq 4$ and $S$ is jointly
closed. By the size-2 and size-3 analyses above and the strict
chain bottom $S_{1} = \{0\}$, the smallest possible $S$ with
$|S| \geq 4$ that is jointly closed is $\Cfour$. To establish
$\Cfour \subseteq S$ for any $|S| \geq 4$ jointly closed: if any
$j \in \{1,\dots,7\} \setminus \{7\}$ lies in $S$, then $B(j,j) =
j+1$ must lie in $S$, so the $B$-orbit on the diagonal forces
$\{j, j+1, j+2, \dots, 7, 8, 9, 0\}$ to lie in $S$ (with
$B(7,7)=8, B(8,8)=7, B(9,9)=0$ providing the closure of the
$\{0,7,8,9\}$ portion); a similar diagonal argument on $\{8, 9\}$
together with $T$'s row-$7$ structure pins down the four
elements of $\Cfour$. Hence any $S$ with $|S| \geq 4$ contains
$\Cfour$, and the listed $S_k$ for $k \geq 4$ are obtained by
adjoining specific operators in the order of the chain.

Direct enumeration over the $\binom{10}{k}$ candidates at each
size (totalling $210+252+210+120+45+10=847$ subsets across sizes
$4$–$9$, plus $\binom{10}{10}=1$ at size $10$) confirms that each
non-listed candidate fails $T$- or $B$-closure on at least one
cell, completing the uniqueness verification. The verification
script (\S\ref{sec:verify}) records the failing cell for each
non-jointly-closed candidate.

The full enumeration of all $1023$ non-empty subsets of $\Zten$,
with the distinguishing failing cell recorded for each
non-jointly-closed subset, is performed by the verification
script \texttt{4core\_verification.py} (\S\ref{sec:verify}).
\end{proof}

\begin{remark}[Structural reading of the chain]\label{rem:sigma-walk}
The chain has a clean structural description in terms of a specific
permutation $\sigma:\Zten\to\Zten$ defined by
\begin{equation}\label{eq:sigma}
\sigma=(0)(3)(8)(9)(1\;7\;6\;5\;4\;2)
\quad\text{in cycle notation.}
\end{equation}
The operators $\{0,3,8,9\}$ are $\sigma$-fixed; the operators
$\{1,2,4,5,6,7\}$ form a single $6$-cycle. We do not derive
$\sigma$ from the tables $T$ and $B$ in the present paper; it is
the canonical operator-cycle permutation arising from the
operator-substrate construction described in
\cite{SandersGish2026Sigma} (where the same $\sigma$ governs the
$\sigma$-rate decomposition of the family $\CL_N$ as $N$ varies).
The structural reading below is then the observation that this
specific $\sigma$ also organizes the chain of
Theorem~\ref{thm:chain}. The chain proceeds:
\begin{itemize}\setlength\itemsep{0pt}
\item Step $1\to 4$: adds $\{7,8,9\}$. Of these, $7$ is in the
$6$-cycle and $\{8,9\}$ are $\sigma$-fixed.
\item Steps $4\to 5$, $5\to 6$, $6\to 7$: each adds the next
element of the $\sigma$-forward orbit of $7$, namely
$\sigma(7)=6$, $\sigma^{2}(7)=5$, $\sigma^{3}(7)=4$.
\item Step $7\to 8$: adds $3$, the remaining $\sigma$-fixed
operator. The $\sigma$-orbit walk pauses for one step:
$\sigma^{4}(7)=2$ would be the next $\sigma$-orbit element, but
adding $2$ alone forces $3=B(2,2)$ into the closure as well,
producing a size jump of $2$. Adding $3$ alone closes more
cheaply.
\item Steps $8\to 9$, $9\to 10$: completes the $\sigma$-orbit
walk, adding $\sigma^{4}(7)=2$ at size $9$ and
$\sigma^{5}(7)=1$ at size $10$.
\end{itemize}
The $\sigma$-fixed operators $\{0,3,8,9\}$ enter at three
positions: $0$ at the bottom (size $1$), $\{8,9\}$ in the
size-$1\to 4$ jump, and $3$ at the size-$7\to 8$ bridge. The
$\sigma$-cycle $\{1,7,6,5,4,2\}$ enters in $\sigma$-forward order,
with the size-$8$ bridge interrupting between
$4=\sigma^{3}(7)$ and $2=\sigma^{4}(7)$.
\end{remark}

\begin{corollary}[Forbidden sizes]\label{cor:forbidden}
No subset of $\Zten$ of size $2$ or size $3$ is jointly closed under
$T$ and $B$.
\end{corollary}

\begin{proof}
Part~(b) of the proof of Theorem~\ref{thm:chain}.
\end{proof}

%% -------- section 4: 4-core --------
\section{The 4-core is jointly closed}\label{sec:4core}

The minimal non-trivial element of the chain in
Theorem~\ref{thm:chain} is the four-element subset
$\Cfour=\{0,7,8,9\}$.

\begin{corollary}[4-core closure]\label{cor:4core}
$\Cfour$ is jointly closed under $T$ and $B$. Explicitly,
$T(\Cfour\times\Cfour)\subseteq\{0,7\}\subset\Cfour$ and
$B(\Cfour\times\Cfour)=\Cfour$.
\end{corollary}

\begin{proof}
Part~(a) of the proof of Theorem~\ref{thm:chain}, displayed
explicitly via the $4\times 4$ restrictions.
\end{proof}

\begin{remark}\label{rem:rank-disparity}
The image of $T\big|_{\Cfour\times\Cfour}$ is the proper subset
$\{0,7\}\subset\Cfour$. The image of $B\big|_{\Cfour\times\Cfour}$
is the full $\Cfour$. The two operations therefore differ
markedly in rank on $\Cfour$. The two restrictions agree on
exactly $4$ of the $16$ cells: $(0,0)\to 0$, $(0,7)\to 7$,
$(7,0)\to 7$, and $(8,8)\to 7$. They disagree on the remaining
$12$ cells. This makes the normalizer identity of
Theorem~\ref{thm:norm} striking, as we now show.
\end{remark}

%% -------- section 5: normalizer --------
\section{The normalizer identity on the 4-core}\label{sec:norm}

For $p\in\R^{10}$ supported on $\Cfour$, write
$(v,h,\beta,r)=(p_{0},p_{7},p_{8},p_{9})$.

\begin{theorem}[Normalizer coincidence]\label{thm:norm}
For every $p\in\R^{10}$ supported on $\Cfour$,
\[
Z_{T}(p)\;=\;Z_{B}(p)\;=\;(v+h+\beta+r)^{2}.
\]
\end{theorem}

\begin{proof}
We compute the four 4-core-supported coordinates of $\Tfuse(p)$
and $\Bfuse(p)$ from Definition~\ref{def:fuse} and the restricted
tables of Corollary~\ref{cor:4core}.

For $T$: cells of $T\big|_{\Cfour\times\Cfour}$ producing $0$ are
$(0,0),(0,8),(8,0),(0,9),(9,0)$ ($5$ cells). Cells producing $7$
are the remaining $11$ cells. Tabulating fuse coordinates in
$(v,h,\beta,r)$:
\begin{align*}
\Tfuse(p)_{0}&=v^{2}+2v\beta+2vr,\\
\Tfuse(p)_{7}&=2vh+h^{2}+2h\beta+2hr+\beta^{2}+2\beta r+r^{2},\\
\Tfuse(p)_{8}&=0,\qquad \Tfuse(p)_{9}=0.
\end{align*}
Summing over $\Cfour$ (and noting $\Tfuse(p)_{c}=0$ for
$c\notin\Cfour$ when $p$ is $\Cfour$-supported, since $T(i,j)\in\Cfour$
for $i,j\in\Cfour$):
\[
Z_{T}(p)\;=\;v^{2}+h^{2}+\beta^{2}+r^{2}+2(vh+v\beta+vr+h\beta+hr+\beta r)\;=\;(v+h+\beta+r)^{2}.
\]

For $B$: cells of $B\big|_{\Cfour\times\Cfour}$ producing $0$ are
$(0,0),(7,9),(9,7),(9,9)$. Cells producing $7$ are
$(0,7),(7,0),(8,8)$. Cells producing $8$ are
$(0,8),(8,0),(7,7),(8,9),(9,8)$. Cells producing $9$ are
$(0,9),(9,0),(7,8),(8,7)$. Tabulating fuse coordinates:
\begin{align*}
\Bfuse(p)_{0}&=v^{2}+2hr+r^{2},\\
\Bfuse(p)_{7}&=2vh+\beta^{2},\\
\Bfuse(p)_{8}&=2v\beta+h^{2}+2\beta r,\\
\Bfuse(p)_{9}&=2vr+2h\beta.
\end{align*}
Summing:
\[
Z_{B}(p)\;=\;v^{2}+h^{2}+\beta^{2}+r^{2}+2(vh+v\beta+vr+h\beta+hr+\beta r)\;=\;(v+h+\beta+r)^{2}.
\]
Both equal $(v+h+\beta+r)^{2}$.
\end{proof}

\begin{corollary}[Normalizer is one on the 4-core simplex]\label{cor:Z=1}
For $p$ in the standard $9$-simplex $\Delta^{9}$ with support in
$\Cfour$, $Z_{T}(p)=Z_{B}(p)=1$.
\end{corollary}

\begin{remark}
The identity $Z_{T}=Z_{B}$ on $\Cfour$ arises by a global
cancellation across the $16$ cells of the restriction, not by
pointwise equality. The two operations agree on only $4$ of the
$16$ cells; the remaining $12$ cells route output values
differently between the four output buckets, in such a way that
the total per-bucket sums coincide. The identity is in fact a
polynomial identity in $(v,h,\beta,r)$, exact and independently
verifiable by symbolic expansion.
\end{remark}

%% -------- section 6: verification --------
\section{Verification}\label{sec:verify}

The two theorems of this paper are independently verified by the
script \texttt{4core\_verification.py}, archived at DOI
\texttt{10.5281/zenodo.18852047}. The script implements:
\begin{enumerate}\setlength\itemsep{0pt}
\item \textbf{Joint-closure enumeration} (\S\ref{sec:chain}):
tests joint closure of all $1023$ non-empty subsets of $\Zten$
and verifies the chain structure of Theorem~\ref{thm:chain},
recording the failing cell for each non-jointly-closed
candidate.
\item \textbf{Normalizer identity} (\S\ref{sec:norm}): symbolically
expands $Z_{T}(p)$ and $Z_{B}(p)$ on $\Cfour$ via \texttt{sympy}
and confirms both equal $(v+h+\beta+r)^{2}$.
\end{enumerate}
The script also implements four further checks supporting the
companion papers in preparation (closed-form attractor at
$\alpha=\tfrac12$; common fixed point across the chain; Galois
$D_{4}$ structure; mixing-weight $\alpha$-sweep); see
\S\ref{sec:forward}. The script runs to completion in
approximately three seconds on a modest desktop CPU. The output
is deterministic.

%% -------- section 7: scope --------
\section{Scope and limits}\label{sec:scope}

The results of this paper are about the specific tables $T$ and
$B$ of Definition~\ref{def:tables}. We do \emph{not} claim:
\begin{itemize}\setlength\itemsep{0pt}
\item that the chain structure of Theorem~\ref{thm:chain} is
shared by other natural pairs of binary operations on $\Zten$;
the chain is a property of the specific $(T, B)$ given;
\item that the normalizer identity of Theorem~\ref{thm:norm}
arises from a categorical or group-theoretic universal
construction; the identity is a polynomial identity in four
coordinates and the cancellation is verified cell-by-cell;
\item that the dynamics of any associated convex-combination
iteration on the simplex is treated here; the present paper
contains no dynamical content. Dynamical results for the pair
$(T, B)$, including a closed-form fixed point of the
convex-combination iteration at $\alpha = \tfrac{1}{2}$ with
explicit algebraic ratios, are the subject of a companion paper
in preparation (see \S\ref{sec:forward}).
\end{itemize}

We \emph{do} claim that the two theorem-level statements of
\S\ref{sec:chain} and \S\ref{sec:norm}, the corollaries of
\S\ref{sec:4core}, and the structural reading of
Remark~\ref{rem:sigma-walk}, are true of the specific $T$ and
$B$ given.

%% -------- section 8: forward directions --------
\section{Forward directions}\label{sec:forward}

The pair $(T, B)$ admits further structural and dynamical
phenomena, treated in companion work in preparation. The
convex-combination iteration $F_{\alpha} = \alpha \widehat{T} +
(1-\alpha)\widehat{B}$ on the $9$-simplex $\Delta^{9}$ has, at
$\alpha = \tfrac{1}{2}$, a unique non-degenerate
$\Cfour$-supported fixed point $p^{*}$ with $p^{*}_{7}/p^{*}_{8}
= 1+\sqrt{3}$ exactly; the ratio $p^{*}_{9}/p^{*}_{8}$ is the
unique positive root of an irreducible integer quartic with
Galois group $D_{4}$ over $\Q$, generating the quartic number
field LMFDB \textnumero 4.2.10224.1. Whether $\alpha =
\tfrac{1}{2}$ is the \emph{unique} rational mixing weight in
$(0,1)\cap\Q$ at which the runtime ratio admits a
small-coefficient algebraic relation is an open question, with
PSLQ evidence (at fifty-digit precision against integer
quadratics with coefficients bounded by $20$) confirming the
algebraic relation only at $\alpha = \tfrac{1}{2}$ among the
sample $\{0, \tfrac{1}{4}, \tfrac{1}{2}, \tfrac{3}{4}, 1\}$. We
also note that the chain of Theorem~\ref{thm:chain} and the
normalizer identity of Theorem~\ref{thm:norm} persist when the
tables are lifted to $\mathbb{F}_{p}^{4}$ for
$p \in \{2, 3, 5, 7, 11, 13\}$.

The companion paper~\cite{SandersGish2026Sigma} provides the
operator-substrate construction from which the pair $(T, B)$ at
$N = 10$ is extracted as a specific instance.

%% -------- references --------
\begin{thebibliography}{99}

\bibitem{DummitFoote2004}
D.~S.~Dummit and R.~M.~Foote.
\newblock \emph{Abstract Algebra}, 3rd edition.
\newblock Wiley, Hoboken, NJ, 2004.

\bibitem{DrapalKepka1985}
A.~Dr\'{a}pal and T.~Kepka.
\newblock Group distances of Latin squares.
\newblock \emph{Commentationes Mathematicae Universitatis Carolinae},
26(2):275--283, 1985.

\bibitem{Kepka1980}
T.~Kepka.
\newblock A note on associative triples of elements in cancellation groupoids.
\newblock \emph{Commentationes Mathematicae Universitatis Carolinae},
21(3):479--487, 1980.

\bibitem{SandersGish2026Sigma}
B.~R.~Sanders and M.~Gish.
\newblock Non-Associativity Decay in Binary Composition Tables over
$\mathbb{Z}/N\mathbb{Z}$.
\newblock Preprint, 2026.
\newblock Companion paper; Zenodo DOI 10.5281/zenodo.18852047.

\end{thebibliography}

\end{document}
